% preamble-exam-zh.tex — 极简讲解页共享 preamble
% 目标：复刻「题目独立、提示轻框、路线箭头、主体双栏、答案轻收束」的讲义风格。

\documentclass{exam-zh}

% ---------- 基础配置 ----------
\examsetup{
  page/size = a4paper,
  page/show-head = false,
  page/show-foot = false,
  sealline/show = false,
  font = none,
  math-font = none,
}

% ---------- 包 ----------
\usepackage{geometry}
\geometry{top=10mm,bottom=14mm,left=12mm,right=10mm}
\AtBeginDocument{\newgeometry{top=10mm,bottom=14mm,left=12mm,right=10mm}}

\usepackage{xcolor}
\usepackage{needspace}
\usepackage{enumitem}
\usepackage{array}
\usepackage{tabularx}
\usepackage{tcolorbox}
\tcbuselibrary{skins,breakable}
\usepackage{paracol}
\usepackage{graphicx}

% ---------- 打印/屏幕颜色切换 ----------
% 用法：\documentclass[monochrome]{exam-zh} 可降级为灰度。
% 注意：不要使用 \if@xxx 放在 makeatother 外部；这里改成普通条件名。
\newif\ifeduprintcolor
\eduprintcolortrue
\makeatletter
\@ifclasswith{exam-zh}{monochrome}{\eduprintcolorfalse}{}
\makeatother

% ---------- 语义颜色 ----------
\ifeduprintcolor
  \definecolor{edu-blue}{RGB}{31,62,126}
  \definecolor{edu-blue-soft}{RGB}{232,238,250}
  \definecolor{edu-blue-bg}{RGB}{248,250,254}
  \definecolor{edu-ink}{RGB}{30,35,45}
  \definecolor{edu-muted}{RGB}{118,124,136}
  \definecolor{edu-line}{RGB}{209,218,235}
  \definecolor{edu-soft-bg}{RGB}{250,251,253}
  \definecolor{edu-red}{RGB}{168,58,58}
  \definecolor{edu-red-bg}{RGB}{255,247,245}
\else
  \definecolor{edu-blue}{gray}{0.15}
  \definecolor{edu-blue-soft}{gray}{0.90}
  \definecolor{edu-blue-bg}{gray}{0.98}
  \definecolor{edu-ink}{gray}{0.10}
  \definecolor{edu-muted}{gray}{0.45}
  \definecolor{edu-line}{gray}{0.82}
  \definecolor{edu-soft-bg}{gray}{0.97}
  \definecolor{edu-red}{gray}{0.28}
  \definecolor{edu-red-bg}{gray}{0.97}
\fi
\colorlet{edu-diagram-muted}{edu-muted}

% ---------- 全局排版 ----------
\setlength{\parindent}{0pt}
\setlength{\parskip}{0pt}
\linespread{1.16}
\renewcommand{\arraystretch}{1.15}

% ctex 通常提供 \kaishu；这里用它做教师旁白感。
\newcommand{\edukai}{\kaishu}

% ---------- 顶部元信息 ----------
\newcommand{\eduTopBar}[2]{%
  \noindent
  \begin{tabularx}{\linewidth}{@{}Xr@{}}
    {\color{edu-blue}\bfseries #1} &
    {\color{edu-blue}\bfseries #2}
  \end{tabularx}
  \vspace{8mm}
}

% ---------- 轻标题体系 ----------
% 只有真正的大结构才用它；不要给提示/路线滥用标题。
\newcommand{\eduSectionTitle}[1]{%
  \needspace{5\baselineskip}%
  \vspace{2mm}
  {\color{edu-blue}\bfseries\large #1}\par
  \vspace{3mm}
}

\newcommand{\eduSubTitle}[1]{%
  \needspace{4\baselineskip}%
  \vspace{1mm}
  {\color{edu-blue}\bfseries #1}\par
  \vspace{1mm}
}

% ---------- 小问题干：复用 exam (1)(2)(3) 格式 ----------
\newenvironment{eduSubQuestionItem}[1]{%
  \needspace{4\baselineskip}%
  \vspace{2mm}
  \par\noindent
  \hangindent=8mm\hangafter=1
  {\color{edu-blue}\bfseries #1}\enspace
  \begingroup
  \color{edu-ink}
}{%
  \endgroup
  \par
  \vspace{3mm}
}

% ---------- 辅助信息条（解题流程/关键想法/思考）楷体 ----------
\newenvironment{eduAuxItem}[1]{%
  \needspace{4\baselineskip}%
  \vspace{2mm}
  \par\noindent
  \hangindent=8mm\hangafter=1
  {\color{edu-blue}\edukai $\triangleright$~#1}\enspace
  \begingroup
  \color{edu-ink}
  \edukai
}{%
  \endgroup
  \par
  \vspace{4mm}
}

\newcommand{\eduColumnTitle}[2]{%
  {\color{edu-blue}\bfseries #1\hspace{1.5mm}#2}\par
  \vspace{2.5mm}
}

% ---------- 图标：不用外部 icon 字体，保证可移植 ----------
\newcommand{\eduIconBulb}{\(\circ\)}
\newcommand{\eduIconBook}{\(\square\)}
\newcommand{\eduIconCheck}{\(\checkmark\)}
\newcommand{\eduIconPin}{\(\diamond\)}
\newcommand{\eduIconNote}{\(\triangleright\)}

% ---------- 题目区 ----------
% 题目区是页面入口：弱标签 + 一条长蓝线 + 大题干。
\newenvironment{eduProblemBlock}[1][题目]{%
  \needspace{8\baselineskip}
  {\small\color{edu-muted}#1}\par
  \vspace{1.5mm}
  {\color{edu-blue}\hrule height 0.55pt}
  \vspace{4.5mm}
  \begingroup
  \color{edu-ink}
  \large
  \setlength{\parskip}{2.2mm}
  \linespread{1.20}\selectfont
}{%
  \par
  \endgroup
  \vspace{6mm}
}

\newenvironment{eduSubQuestions}{%
  \begin{enumerate}[
    label=(\arabic*),
    leftmargin=8mm,
    itemsep=2.0mm,
    topsep=2.0mm,
    parsep=0pt
  ]
}{%
  \end{enumerate}
}

% ---------- 读题提示：轻框，不做同级标题 ----------
\newtcolorbox{eduReadTipBox}{
  enhanced,
  breakable,
  colback=white,
  colframe=edu-blue!45,
  boxrule=0.45pt,
  arc=1.6mm,
  left=10mm,
  right=4mm,
  top=5mm,
  bottom=5mm,
  before skip=1mm,
  after skip=7mm,
  fontupper=\edukai\normalsize,
  colupper=edu-ink,
  overlay={
    \node[anchor=north west, text=edu-blue] at ([xshift=3.2mm,yshift=-3.2mm]frame.north west)
    {\small\eduIconNote};
  }
}

\newenvironment{eduTipList}{%
  \begin{itemize}[
    label=\textbullet,
    leftmargin=5mm,
    itemsep=1.2mm,
    topsep=0pt,
    parsep=0pt
  ]
}{%
  \end{itemize}
}

% ---------- 解题路线：虚线框 + 文字箭头 ----------
\newenvironment{eduRouteFlow}{%
  \needspace{4\baselineskip}%
  \vspace{1mm}
  \begin{center}
  \normalsize
}{%
  \end{center}
  \vspace{7mm}
}
\newcommand{\eduRouteStep}[1]{%
  {\color{edu-ink}#1}%
}
\newcommand{\eduRouteArrow}{%
  \;$\boldsymbol{\to}$\;%
}

% ---------- 主体讲解：使用 paracol 实现跨页自适应双栏 ----------
\newenvironment{eduDualExplain}{%
  \needspace{6\baselineskip}% 防止标题孤行
  \columnratio{0.26, 0.74}%
  \setlength{\columnsep}{8mm}%
  \columnseprule=0.45pt%
  \colseprulecolor{edu-blue}%
  \begin{paracol}{2}% 开启双栏
}{%
  \end{paracol}% 结束双栏
  \vspace{5mm}
}

\newenvironment{eduExplainLeft}{%
  \setlength{\linewidth}{\columnwidth}%
  \hsize=\columnwidth
  \normalsize\color{edu-ink}%
}{%
  \switchcolumn
}

\newcommand{\eduColumnDivider}{}

\newenvironment{eduExplainRight}{%
  \setlength{\linewidth}{\columnwidth}%
  \hsize=\columnwidth
  \normalsize\color{edu-ink}%
}{%
}

\newtcolorbox{eduSideHintBox}[1]{
  enhanced,
  breakable,
  colback=edu-blue-bg,
  colframe=edu-line,
  boxrule=0.35pt,
  borderline west={1.5pt}{0pt}{edu-blue},
  arc=0.8mm,
  left=2.4mm,
  right=2.4mm,
  top=2mm,
  bottom=2mm,
  before skip=1mm,
  after skip=1.5mm,
  title={\color{edu-blue}\bfseries\small #1},
  coltitle=edu-blue,
  attach title to upper={\par\vspace{0.7mm}},
  fontupper=\small
}

\newtcolorbox{eduSideMistakeBox}[1]{
  enhanced,
  breakable,
  colback=edu-red-bg,
  colframe=edu-red!35,
  boxrule=0.35pt,
  borderline west={1.5pt}{0pt}{edu-red},
  arc=0.8mm,
  left=2.4mm,
  right=2.4mm,
  top=2mm,
  bottom=2mm,
  before skip=1mm,
  after skip=1.5mm,
  title={\color{edu-red}\bfseries\small #1},
  coltitle=edu-red,
  attach title to upper={\par\vspace{0.7mm}},
  fontupper=\small
}

\newtcolorbox{eduSideNoteBox}[1]{
  enhanced,
  breakable,
  colback=edu-soft-bg,
  colframe=edu-line,
  boxrule=0.35pt,
  borderline west={1.5pt}{0pt}{edu-muted},
  arc=0.8mm,
  left=2.4mm,
  right=2.4mm,
  top=2mm,
  bottom=2mm,
  before skip=1mm,
  after skip=1.5mm,
  title={\color{edu-muted}\bfseries\small #1},
  coltitle=edu-muted,
  attach title to upper={\par\vspace{0.7mm}},
  fontupper=\small
}

\newcommand{\eduSideHint}[2]{%
  \begin{eduSideHintBox}{#1}#2\end{eduSideHintBox}%
}
\newcommand{\eduSideMistake}[2]{%
  \begin{eduSideMistakeBox}{#1}#2\end{eduSideMistakeBox}%
}
\newcommand{\eduSideNote}[2]{%
  \begin{eduSideNoteBox}{#1}#2\end{eduSideNoteBox}%
}

\newcommand{\eduExplainStep}[3]{%
  \needspace{3\baselineskip}%
  \vspace{1.2mm}%
  {\color{edu-blue}\bfseries step#1.\enspace #2}\par
  \vspace{0.8mm}%
  \begingroup
  \color{edu-ink}#3\par
  \endgroup
  \vspace{2.2mm}%
}

\newenvironment{eduNumberList}{%
  \begin{enumerate}[
    label=\arabic*.,
    leftmargin=6mm,
    itemsep=4.8mm,
    topsep=0pt,
    parsep=0pt
  ]
}{%
  \end{enumerate}
}

\newenvironment{eduBulletList}{%
  \begin{itemize}[
    label=\textbullet,
    leftmargin=5mm,
    itemsep=1.2mm,
    topsep=0pt,
    parsep=0pt
  ]
}{%
  \end{itemize}
}

% ---------- 底部双栏：答案 + 方法提醒 ----------
\newenvironment{eduBottomGrid}{%
  \columnratio{0.32, 0.68}%
  \setlength{\columnsep}{11mm}%
  \columnseprule=0pt%
  \begin{paracol}{2}%
}{%
  \end{paracol}%
}

\newenvironment{eduSummaryLeft}{%
}{%
}

\newcommand{\eduSummaryDivider}{%
  \switchcolumn
}

\newenvironment{eduSummaryRight}{%
}{%
}

% ---------- 答案/方法提醒：轻量收束 ----------
\newtcolorbox{eduSummaryBlock}[2]{
  enhanced,
  breakable,
  colback=edu-blue-bg,
  colframe=edu-line,
  boxrule=0.45pt,
  arc=1.2mm,
  left=3.5mm,
  right=3.5mm,
  top=3mm,
  bottom=3mm,
  before skip=1mm,
  after skip=2.5mm,
  title={\color{edu-blue}\bfseries #1\hspace{1.5mm}#2},
  coltitle=edu-blue,
  attach title to upper={\par\vspace{1.5mm}},
  fontupper=\normalsize
}

% ---------- 例题单栏解答卡片 ----------
\newtcolorbox{eduSolutionBox}[1]{
  enhanced,
  breakable,
  colback=white,
  colframe=edu-blue!40,
  boxrule=0.55pt,
  arc=1.2mm,
  left=4mm,
  right=4mm,
  top=3.5mm,
  bottom=3.5mm,
  before skip=2mm,
  after skip=3mm,
  title={\color{edu-blue}\bfseries \eduIconBook\hspace{1.5mm}#1},
  coltitle=edu-blue,
  attach title to upper={\par\vspace{1.5mm}},
  fontupper=\normalsize
}

% ---------- 变式训练：完整题目 + 学生留白 ----------
\newtcolorbox{eduVariationTraining}[1]{
  enhanced,
  breakable,
  colback=white,
  colframe=edu-blue!45,
  boxrule=0.5pt,
  borderline west={1.8pt}{0pt}{edu-blue},
  arc=1.2mm,
  left=4mm,
  right=4mm,
  top=3.2mm,
  bottom=3.2mm,
  before skip=2mm,
  after skip=4mm,
  title={\color{edu-blue}\bfseries #1},
  coltitle=edu-blue,
  attach title to upper={\par\vspace{1.8mm}},
  fontupper=\normalsize,
  colupper=edu-ink
}

\newenvironment{eduPlainList}{%
  \begin{enumerate}[
    label=(\arabic*),
    leftmargin=7mm,
    itemsep=1.2mm,
    topsep=0pt,
    parsep=0pt
  ]
}{%
  \end{enumerate}
}

% ---------- 兼容旧 block 的轻量盒子 ----------
\newtcolorbox{eduSoftBox}{
  enhanced,
  breakable,
  colback=edu-soft-bg,
  colframe=edu-line,
  boxrule=0.35pt,
  arc=1mm,
  left=3mm,
  right=3mm,
  top=2mm,
  bottom=2mm,
  before skip=1mm,
  after skip=2mm,
  fontupper=\small
}

\newtcolorbox{eduMistakeBox}{
  enhanced,
  breakable,
  colback=edu-red-bg,
  colframe=edu-red!40,
  boxrule=0.35pt,
  arc=1mm,
  left=3mm,
  right=3mm,
  top=2.5mm,
  bottom=2.5mm,
  before skip=1mm,
  after skip=2mm,
  fontupper=\normalsize
}

\newtcolorbox{eduLegacySolution}{
  enhanced,
  breakable,
  colback=edu-soft-bg,
  colframe=edu-line,
  boxrule=0.35pt,
  arc=1mm,
  left=3mm,
  right=3mm,
  top=2mm,
  bottom=2mm,
  before skip=-2mm,
  after skip=4mm,
  fontupper=\small
}

\newtcolorbox{eduTeacherNote}{
  enhanced,
  breakable,
  colback=edu-soft-bg,
  colframe=edu-line,
  boxrule=0pt,
  borderline west={1.5pt}{0pt}{edu-muted},
  left=2.5mm,
  right=2.5mm,
  top=2mm,
  bottom=2mm,
  before skip=1mm,
  after skip=2mm,
  fontupper=\footnotesize,
  colupper=edu-muted
}

% ---------- 答题区 ----------
\newcommand{\answerarea}[1][40mm]{%
  \vspace{3mm}%
  \begin{tcolorbox}[
    enhanced,
    breakable,
    colback=white,
    colframe=edu-line,
    boxrule=0.55pt,
    arc=0.8mm,
    height=#1,
    valign=top,
  ]
  \end{tcolorbox}%
}

\newcommand{\diagramcolfigure}[3]{%
  \begin{minipage}[t]{#2}
    \vspace{0pt}\centering
    \includegraphics[width=\linewidth]{\detokenize{#1}}
    \if\relax\detokenize{#3}\relax\else
      \par\vspace{1mm}{\scriptsize\color{edu-diagram-muted}#3}
    \fi
  \end{minipage}%
}

\newcommand{\diagramrowitem}[4]{%
  \begin{minipage}[t]{#2}
    \vspace{0pt}\centering
    \includegraphics[width=\linewidth]{\detokenize{#1}}
    \if\relax\detokenize{#3}\relax\else
      \par\vspace{1mm}{\scriptsize\bfseries #3}
    \fi
    \if\relax\detokenize{#4}\relax\else
      \par\vspace{0.5mm}{\scriptsize\color{edu-diagram-muted}#4}
    \fi
  \end{minipage}%
}

\newcommand{\answerareawithdiagramcol}[4]{%
  \vspace{3mm}%
  \begin{tcolorbox}[
    enhanced,
    breakable,
    colback=white,
    colframe=edu-line,
    boxrule=0.55pt,
    arc=0.8mm,
    height=#1,
    valign=top,
  ]
  \begin{minipage}[t]{\dimexpr\linewidth-#3-6mm\relax}
    \vspace{0pt}
  \end{minipage}\hfill
  \diagramcolfigure{#2}{#3}{#4}
  \end{tcolorbox}%
}

\newcommand{\answerareawithdiagram}[4]{%
  \answerareawithdiagramcol{#1}{#2}{#3}{#4}%
}

\examsetup{
  question/show-answer = true,
  fillin/show-answer = true,
  solution/show-solution = show-stay,
}

\begin{document}

\eduSectionTitle{知识点 4：待定系数法求二次函数解析式——根据信息选式}


\begin{eduSolutionBox}{选式决策表：什么信息配什么式}
  \begingroup
  \setlength{\parskip}{2.5mm}
\textbf{核心：} 二次函数有 3 个自由度（$a,b,c$），需 3 个独立条件唯一确定。\par
\begin{tabular}{p{0.34\linewidth}p{0.34\linewidth}p{0.22\linewidth}}
题目给的信息 & 应设的式 & 待定参数\\
三个普通点 & 一般式 $y=ax^2+bx+c$ & $a,b,c$\\
顶点 + 一个点 & 顶点式 $y=a(x-h)^2+k$ & $a$（$h,k$ 已知）\\
对称轴 + 两点 & 顶点式（$h$ 已知）& $a,k$\\
两个 x 轴交点 + 一个点 & 交点式 $y=a(x-x_1)(x-x_2)$ & $a$\\
y 轴交点 $(0,c)$ & 一般式（$c$ 已知）& $a,b$\\
\end{tabular}
\par
\textbf{选式口诀：} 有顶点设顶点式，有 x 轴交点设交点式，其余设一般式。选对式，方程最少。\par
  \endgroup
\end{eduSolutionBox}




\begin{eduProblemBlock}[例题 1（三个普通点 → 设一般式）]
求过三点 $(0,3),(1,0),(2,1)$ 的二次函数解析式。

\end{eduProblemBlock}





\begin{eduAuxItem}{思路导航}
选式：三个普通点 → 设一般式$\;\to\;$由 $(0,3)$ 定 $c$$\;\to\;$由另两点列二元方程组求 $a,b$$\;\to\;$解方程组$\;\to\;$写出解析式\end{eduAuxItem}




\begin{eduSubQuestionItem}{例题 1}
求过 $(0,3),(1,0),(2,1)$ 的二次函数解析式。\end{eduSubQuestionItem}

\begin{eduDualExplain}
  \begin{eduExplainLeft}
    \eduColumnTitle{\eduIconBulb}{小贴士}
        \eduSideHint{为什么先代 $(0,c)$}{$x=0$ 代入直接得 $c$，最省事；剩下只需解二元方程组。}
        \eduSideMistake{消元别算错}{$2a+b=-1$ 减去 $a+b=-3$ 得 $a=2$，不是 $a=-2$；注意符号。}
  \end{eduExplainLeft}

  \begin{eduExplainRight}
    \eduColumnTitle{\eduIconBook}{解答}
      \eduExplainStep{1}{选式：三个普通点 → 设一般式}{三点都不是顶点、也不是 x 轴成对交点，设 $y=ax^2+bx+c$。}
      \eduExplainStep{2}{由 $(0,3)$ 定 $c$}{代入 $(0,3)$：$c=3$。}
      \eduExplainStep{3}{由另两点列二元方程组求 $a,b$}{代入 $(1,0)$：$a+b=-3$；代入 $(2,1)$：$2a+b=-1$。}
      \eduExplainStep{4}{解方程组}{两式相减得 $a=2$，代回得 $b=-5$。}
      \eduExplainStep{5}{写出解析式}{$y=2x^2-5x+3$。}
    \vspace{1mm}
    {\color{edu-blue}\bfseries\small 验算}\par\vspace{1mm}
    \begin{eduBulletList}
      \item $(1,0)$：$2-5+3=0$ ✓；$(2,1)$：$8-10+3=1$ ✓。    \end{eduBulletList}
  \end{eduExplainRight}
\end{eduDualExplain}




\begin{eduProblemBlock}[例题 2（顶点 + 一个点 → 设顶点式）]
已知二次函数的顶点为 $(1,-2)$，且图像经过点 $(2,0)$，求解析式。

\end{eduProblemBlock}





\begin{eduAuxItem}{思路导航}
选式：有顶点 → 设顶点式$\;\to\;$代入 $(2,0)$ 定 $a$$\;\to\;$写出解析式\end{eduAuxItem}




\begin{eduSubQuestionItem}{例题 2}
顶点 $(1,-2)$，过 $(2,0)$，求解析式。\end{eduSubQuestionItem}

\begin{eduDualExplain}
  \begin{eduExplainLeft}
    \eduColumnTitle{\eduIconBulb}{小贴士}
        \eduSideMistake{不要设一般式}{若设 $y=ax^2+bx+c$，顶点条件要翻译成 $-\frac{b}{2a}=1$ 和 $\frac{4ac-b^2}{4a}=-2$，再加一个点，列三元方程组，麻烦得多。}
        \eduSideHint{顶点式的好处}{顶点已知时，$h,k$ 直接代入，只剩 1 个未知数 $a$，只需 1 个方程。}
  \end{eduExplainLeft}

  \begin{eduExplainRight}
    \eduColumnTitle{\eduIconBook}{解答}
      \eduExplainStep{1}{选式：有顶点 → 设顶点式}{顶点 $(1,-2)$，设 $y=a(x-1)^2-2$，$h=1,k=-2$ 已知，只需定 $a$。}
      \eduExplainStep{2}{代入 $(2,0)$ 定 $a$}{代入 $(2,0)$：$a(2-1)^2-2=0\Rightarrow a-2=0\Rightarrow a=2$。}
      \eduExplainStep{3}{写出解析式}{$y=2(x-1)^2-2$。}
    \vspace{1mm}
    {\color{edu-blue}\bfseries\small 验算}\par\vspace{1mm}
    \begin{eduBulletList}
      \item 展开 $2(x-1)^2-2=2x^2-4x$，顶点 $(1,-2)$：$2-4=-2$ ✓；$(2,0)$：$8-8=0$ ✓。    \end{eduBulletList}
  \end{eduExplainRight}
\end{eduDualExplain}




\begin{eduProblemBlock}[例题 3（两个 x 轴交点 + 一个点 → 设交点式）]
已知二次函数图像与 $x$ 轴交于 $(-1,0),(3,0)$ 两点，且经过点 $(0,-3)$，求解析式。

\end{eduProblemBlock}





\begin{eduAuxItem}{思路导航}
选式：有 x 轴交点 → 设交点式$\;\to\;$代入 $(0,-3)$ 定 $a$$\;\to\;$写出解析式\end{eduAuxItem}




\begin{eduSubQuestionItem}{例题 3}
与 x 轴交于 $(-1,0),(3,0)$，过 $(0,-3)$，求解析式。\end{eduSubQuestionItem}

\begin{eduDualExplain}
  \begin{eduExplainLeft}
    \eduColumnTitle{\eduIconBulb}{小贴士}
        \eduSideMistake{交点式符号}{根是 $-1$ 和 $3$，交点式写成 $a(x-(-1))(x-3)=a(x+1)(x-3)$；根为负时括号里是加。}
        \eduSideHint{为什么交点式省事}{两个根已知，直接消掉 $b,c$（或 $h,k$），只剩 1 个未知数 $a$。}
  \end{eduExplainLeft}

  \begin{eduExplainRight}
    \eduColumnTitle{\eduIconBook}{解答}
      \eduExplainStep{1}{选式：有 x 轴交点 → 设交点式}{两根 $x_1=-1,x_2=3$，设 $y=a(x+1)(x-3)$，只需定 $a$。}
      \eduExplainStep{2}{代入 $(0,-3)$ 定 $a$}{代入 $(0,-3)$：$a(0+1)(0-3)=-3a=-3\Rightarrow a=1$。}
      \eduExplainStep{3}{写出解析式}{$y=(x+1)(x-3)=x^2-2x-3$。}
    \vspace{1mm}
    {\color{edu-blue}\bfseries\small 验算}\par\vspace{1mm}
    \begin{eduBulletList}
      \item $(-1,0)$：$1+2-3=0$ ✓；$(3,0)$：$9-6-3=0$ ✓；$(0,-3)$：$0-0-3=-3$ ✓。    \end{eduBulletList}
  \end{eduExplainRight}
\end{eduDualExplain}




\begin{eduMistakeBox}
\eduSubTitle{易错提醒}\begin{itemize}
  \item \textbf{一律设一般式}：给了顶点或 x 轴交点仍设一般式，导致列三元方程组，计算量爆炸。
  \item \textbf{交点式根的符号}：根为 $-1$ 时括号写 $x+1$（即 $x-(-1)$），不是 $x-1$。
  \item \textbf{顶点 + 对称轴当两个条件}：顶点横坐标就是对称轴，二者是同一条信息，不算 2 个独立条件。
  \item \textbf{忘验证}：求出解析式后必须回代所有给定点检查。
\end{itemize}
\end{eduMistakeBox}




\begin{eduSummaryBlock}{\eduIconPin}{方法提醒}
\begin{eduBulletList}
  \item 3 个独立条件定一个二次函数。  \item 选式口诀：有顶点设顶点式，有 x 轴交点设交点式，其余设一般式。  \item 顶点式/交点式只需定 $a$，一般式需列方程组。  \item 求出后展开/配方回代验证。\end{eduBulletList}
\end{eduSummaryBlock}




\begin{eduBottomGrid}
  \begin{eduSummaryLeft}
    \begin{eduSummaryBlock}{\eduIconCheck}{答案}
    \begin{eduPlainList}
      \item 例 1：$y=2x^2-5x+3$      \item 例 2：$y=2(x-1)^2-2$      \item 例 3：$y=(x+1)(x-3)$    \end{eduPlainList}
    \end{eduSummaryBlock}
  \end{eduSummaryLeft}
  \eduSummaryDivider
  \begin{eduSummaryRight}
    \begin{eduSummaryBlock}{\eduIconPin}{一句话}
    \begin{eduBulletList}
      \item 看信息选式，让未知数最少。    \end{eduBulletList}
    \end{eduSummaryBlock}
  \end{eduSummaryRight}
\end{eduBottomGrid}




\end{document}
